% Part: first-order-logic
% Chapter: models-theories
% Section: size-of-structures

\documentclass[../../../include/open-logic-section]{subfiles}

\begin{document}

\olfileid{fol}{mat}{siz}

\olsection{\printtoken{P}{structure}-এর আকার প্রকাশ}

\begin{explain}
কোনো ভাষার অ-যৌক্তিক সংকেত ব্যবহার না করেও গঠনগুলির কিছু ধর্ম
প্রকাশ করা যায়। যেমন, এমন !!{sentence}s আছে যা কোনো
!!a{structure}-এ ঠিক তখনই সত্য, যখন !!{structure}-টির !!{domain}-এ
অন্তত, বড়জোর, অথবা ঠিক একটি নির্দিষ্ট সংখ্যক~$n$ !!{element}s আছে।
\end{explain}

\begin{prop}
!!{sentence}
\begin{multline*}
  !A_{\ge n} \ident \lexists[x_1][\lexists[x_2][\dots\lexists[x_n][{}]]]\\
  \begin{aligned}
  (\eq/[x_1][x_2] \land {}
  \eq/[x_1][x_3] \land \eq/[x_1][x_4] \land \dots \land \eq/[x_1][x_n] \land {}\\
\eq/[x_2][x_3] \land \eq/[x_2][x_4] \land \dots \land {} \eq/[x_2][x_n] \land {} \\
\vdots\\
\eq/[x_{n-1}][x_n])
\end{aligned}
\end{multline*}
কোনো !!a{structure}~$\Struct M$-এ সত্য যদি এবং কেবল যদি
$\Domain M$-এ অন্তত~$n$ !!{element}s থাকে। ফলে
$\Sat{M}{\lnot !A_{\ge n+1}}$ যদি এবং কেবল যদি $\Domain M$-এ
বড়জোর~$n$ !!{element}s থাকে।

\end{prop}

\begin{prop}
!!{sentence}
\begin{multline*}
!A_{= n} \ident \lexists[x_1][\lexists[x_2][\dots\lexists[x_n][{}]]] \\
\begin{aligned}
  (\eq/[x_1][x_2] \land {}
  \eq/[x_1][x_3] \land \eq/[x_1][x_4] \land \dots \land \eq/[x_1][x_n] \land {}\\
\eq/[x_2][x_3] \land \eq/[x_2][x_4] \land \dots \land {} \eq/[x_2][x_n] \land {} \\
\vdots\\
\eq/[x_{n-1}][x_n] \land {} \\
\lforall[y][(\eq[y][x_1] \lor \dots \lor \eq[y][x_n])])
\end{aligned}
\end{multline*}
% BN-SRC-108 TARGET-CORRECTION-BEGIN
\par\smallskip\noindent\textnormal{\small\textbf{উৎস-সংশোধন (BN-SRC-108):} হিমায়িত ইংরেজি $!A_{=n}$-এর শেষ সারিতে সার্বিক পরিমাণসূচকের সূত্র-আর্গুমেন্ট পূর্বেই বন্ধ করে তার পরে দুটি গোলবন্ধনী রাখা হয়েছে। বাংলা পাঠে বিচ্ছেদটির সমাপ্তি-গোলবন্ধনী আগে এবং পরিমাণসূচকের সমাপ্তি-বর্গবন্ধনী পরে রেখে অতিরিক্ত গোলবন্ধনীটি বাদ দেওয়া হয়েছে।} % BN-SRC-108 TARGET-CORRECTION-END
কোনো !!a{structure}~$\Struct M$-এ সত্য যদি এবং কেবল যদি
$\Domain M$-এ ঠিক~$n$ !!{element}s থাকে।
\end{prop}

\begin{prop}
কোনো !!{structure} অসীম যদি এবং কেবল যদি সেটি
\[
\{!A_{\ge 1}, !A_{\ge 2}, !A_{\ge 3}, \dots \}.
\]
-এর একটি মডেল।
\end{prop}

এমন কোনো একক বিশুদ্ধ যৌক্তিক বাক্য নেই যা~$\Struct M$-এ ঠিক তখনই সত্য,
যখন~$\Domain M$ অসীম। তবে অ-যৌক্তিক !!{predicate}s-সহ এমন
!!{sentence}s দেওয়া যায় যাদের কেবল অসীম মডেল আছে (যদিও প্রত্যেক অসীম
!!{structure} তাদের মডেল নয়)। কোনো গঠন সসীম হওয়ার ধর্ম এবং কোনো গঠন
!!{nonenumerable} হওয়ার ধর্ম, এমনকি !!{sentence}s-এর একটি অসীম সেট
দিয়েও প্রকাশ করা যায় না। সংহতি ও লোয়েনহাইম--স্কোলেম উপপাদ্য থেকে
এই তথ্যগুলি অনুসৃত হয়।

\end{document}
